\documentclass[12pt,a4paper]{article}

\usepackage{fontspec}
\usepackage{polyglossia}
\setdefaultlanguage{english}
\setmainfont{CMU Serif}
\setsansfont{CMU Sans Serif}
\setmonofont{CMU Typewriter Text}
\usepackage{geometry}
\geometry{top=2cm, bottom=2cm, left=2.5cm, right=2.5cm}
\usepackage{amsmath,amssymb,amsthm}
\usepackage{booktabs,longtable,float,xcolor}
\usepackage{hyperref}
\hypersetup{colorlinks, urlcolor=blue, citecolor=blue, linkcolor=black}
\usepackage{enumitem}

\tolerance=2000
\emergencystretch=3em
\hbadness=10000
\vbadness=10000

\numberwithin{equation}{section}

\title{\textbf{Model MEML~75:}\\[0.3cm]
\large Quark and lepton mass matrices from mixed \\
\(T^2/\mathbb{Z}_3 \times T^2/\mathbb{Z}_4\) geometry \\
with axion-like dark matter and leptogenesis}
\author{D.~V.~Yurov}
\date{2026}

\begin{document}
\maketitle

\begin{abstract}
\noindent
We present model MEML~75, which parametrizes the fermion sector of the
Standard Model through the geometry of mixed six-dimensional orbifolds.
The quark sector is built on \(T^2/\mathbb{Z}_3\) using \(A_4\) modular forms,
while the lepton sector is localized on \(T^2/\mathbb{Z}_4\), requiring the
double cover modular group \(S_4'\). The model has \textbf{19 effective parameters}
fitting \textbf{16 observables} (9 quarks + 5 lepton observables + \(\Omega_a h^2\) + \(\eta_B\)),
yielding \(N_{\rm dof}^{\rm fit} = 16 - 19 = -3\). The absolute neutrino mass scale is fixed by the Type-I See-Saw mechanism via the six-dimensional Planck mass \(M_6\). 
The mixed orbifold geometry is naturally embedded in the framework of \emph{eclectic flavor symmetries}.

\textbf{The main testable consequences:}
axion-like dark matter with
\(f_a = 1.9801\times10^{11}\)~GeV,
\(m_a = 28.407\)~\(\mu\)eV, \(\nu_a = 6.8688\)~GHz, and ALP-photon
coupling
\(g_{a\gamma\gamma} = +4.38\times10^{-15}\)~GeV\(^{-1}\) (DFSZ)
or \(-1.13\times10^{-14}\)~GeV\(^{-1}\) (KSVZ), within reach of future
high-frequency haloscope searches such as ADMX and HAYSTAC~\cite{ADMX2018,HAYSTAC2021}.
The model also predicts a sum of neutrino masses \(\sum m_i \approx 0.068\) eV, an effective Majorana mass \(m_{\beta\beta} \approx 8.7\) meV, and a non-zero Jarlskog invariant \(J_{cp} \approx 0.0156\) (\(\delta_{CP} \approx 28^\circ\) or \(152^\circ\)), all testable in near-future experiments.
Of the 19 parameters, 17 are fitted to the 16 observables (the 15 quark/lepton parameters, plus \(\theta_i\) for \(\Omega_a h^2\) and \(\Delta\) for \(\eta_B\)); the remaining 2 are the UV inputs \(y_\nu\) and \(\mathcal{I}\), counted conservatively as free.
Counting the three additional predictions (\(\delta_{CP}\), \(\sum m_i\), \(m_{\beta\beta}\)) yields \(N_{\rm dof}^{\rm total} = (16+3)-19 = 0\), so the model is exactly determined rather than overdetermined.

\textbf{Quark sector (\(T^2/\mathbb{Z}_3\), \(A_4\)).} A global fit yields
\(\chi^2 < 10^{-10}\) (numerical precision limit) at
\(\tau_u = -0.1577944085 + 2.3851782590\,i\) and
\(\tau_d = +0.3693905617 + 1.5410496057\,i\).

\textbf{Lepton sector (\(T^2/\mathbb{Z}_4\), \(S_4'\)).} Using the weight-2 modular forms of \(S_4'\) and a uniform diagonal shift, the neutrino mass matrix successfully reproduces the PMNS angles and mass-squared differences. The global fit yields \(\chi^2 \approx 0.23\) at \(\tau_l = 1.5398 + 1.5511\,i\), predicting \(\theta_{12}=33.50^\circ\), \(\theta_{13}=8.53^\circ\), \(\theta_{23}=49.00^\circ\).

\textbf{Verdict:} the model is highly successful. The mixed orbifold geometry allows the model to accurately describe both the quark hierarchy and the leptonic mixing, while maintaining multiple testable predictions for ALP dark matter, neutrino masses, CP violation, and the origin of matter.
\end{abstract}

\newpage

\section{Introduction}
Modular models of fermion masses~\cite{Feruglio2017,Kobayashi2019,Novichkov2019,Ding2021}
build the mass hierarchy and mixing from compactification geometry.
Recently, the concept of \emph{eclectic flavor symmetry} has emerged as a powerful framework, combining traditional discrete symmetries with modular symmetries that depend on the compactification moduli~\cite{Nilles2020,Baur2024}. In model MEML~75, we explore a mixed geometric setup: quarks are localized on a \(T^2/\mathbb{Z}_3\) torus with moduli \(\tau_u\) and \(\tau_d\), while leptons are localized on a \(T^2/\mathbb{Z}_4\) torus with modulus \(\tau_l\). This dual-orbifold approach allows us to leverage the best features of two different modular groups: \(A_4\) for the hierarchical quark masses, and \(S_4'\) for the rich leptonic mixing pattern.

\section{Geometry of \(T^2/\mathbb{Z}_3\) and \(T^2/\mathbb{Z}_4\)}
The orbifold \(T^2/\mathbb{Z}_3\) is defined by \(z = x + \tau y\) with \(\tau = e^{i\pi/3}\). Its $Z_3$-action has three fixed points, providing a topological explanation for three generations. The quark mass matrices are built from \(A_4\) modular forms of weights \((2,4,6)\).

The mixed geometry \(T^2/\mathbb{Z}_3 \times T^2/\mathbb{Z}_4\) naturally arises in Type IIB orientifold compactifications with D-branes wrapping different cycles. Specifically, one can consider a \(T^6/(\mathbb{Z}_3 \times \mathbb{Z}_4)\) orbifold, where the six-dimensional torus factorizes into two invariant subtori. The discrete torsion or flux background can select different orbifold groups for quarks and leptons. This places the model within the general framework of eclectic flavor symmetries~\cite{Baur2024}. The lepton sector uses the \(T^2/\mathbb{Z}_4\) orbifold, whose modular symmetry is the double cover group \(S_4'\)~\cite{Novichkov2019,Ding2021}.

\section{Bulk coupling \(\alpha_G\) (postulate)}
\begin{equation}
\alpha_G = \frac{2}{3\pi} = 0.2122065908.
\end{equation}
Geometric motivation: \(2/3\) is the volume fraction of the fundamental domain of the \(\mathbb{Z}_3\)-orbifold; \(1/\pi\) is the normalization of the theta-function period.

\section{Suppression factor \(\varepsilon\) (ansatz)}
\begin{equation}
\varepsilon = \left(\frac{\alpha_G}{4\pi}\right)^2 = \frac{1}{36\pi^4} = 2.851662\times10^{-4}.
\end{equation}
The factor \(4\pi\) in the denominator can be read as the standard two-loop suppression \(1/(16\pi^2)\).

\section{Conditionally derived quantities}
\begin{align}
m_{KK} &= \frac{M_{Pl}^{\rm red}}{\sqrt{\mathcal{I}}\,e^{\pi^2}} = 2.0838\times10^{14}~\text{GeV}, \\
M_6 &= \left(3 (M_{Pl}^{\rm red})^2 m_{KK}^2\right)^{1/4} = 2.9646\times10^{16}~\text{GeV},
\end{align}
at \(\mathcal{I} = 0.3653\).\footnote{The geometric factor \(\mathcal{I} = 0.3653\) and the neutrino Yukawa coupling \(y_\nu = 0.934\) are \emph{not} derived from first principles in this work. To be conservative, we count them as free parameters in the balance of Section~10, giving $N_{\rm eff}=19$ and $N_{\rm dof}^{\rm total}=0$. A future UV derivation of these two inputs would reduce $N_{\rm eff}$ to 17 and restore $N_{\rm dof}^{\rm total}=+2$.}

\section{Axion-like (ALP) sector}
\subsection{ALP parameters and dark matter}
\begin{align}
f_a &= M_{Pl}^{\rm red}\,\varepsilon^2 = 1.9801\times10^{11}~\text{GeV}, \\
m_a &= \frac{(\chi^{1/4})^2}{f_a} = 28.407~\mu\text{eV}, \\
\nu_a &= \frac{m_a c^2}{h} = 6.8688~\text{GHz}.
\end{align}
Here \(\chi^{1/4} \approx 75.0\)~MeV is the fourth root of the topological susceptibility of QCD. The relic density is:
\begin{equation}
\Omega_a h^2 \approx 0.12\,\theta_i^2\, f_{\rm anh}(\theta_i)\left(\frac{f_a}{3\times10^{11}}\right)^{7/6},
\end{equation}
where \(f_{\rm anh}(\theta_i) = [\ln(e/(1-(\theta_i/\pi)^2))]^{7/6}\) is the anharmonicity factor. We treat $\theta_i$ as a free parameter and fit it to the Planck value $\Omega_a h^2 = 0.12$~\cite{Planck2018}, obtaining \(\theta_i \approx 1.17\).

\subsection{ALP-photon coupling and superradiance}
\begin{equation}
g_{a\gamma\gamma} = \frac{\alpha_{\rm EM}}{2\pi f_a}\left(\frac{E}{N} - 1.92\right).
\end{equation}
At \(f_a = 1.9801\times10^{11}\)~GeV: \(g_{a\gamma\gamma} = +4.38\times10^{-15}\)~GeV\(^{-1}\) (DFSZ) or \(-1.13\times10^{-14}\)~GeV\(^{-1}\) (KSVZ). The predicted mass \(m_a = 28.4~\mu\)eV corresponds to \(\nu_a = 6.87\)~GHz, within reach of future high-frequency haloscope searches such as ADMX~\cite{ADMX2018} and HAYSTAC~\cite{HAYSTAC2021}. The ALP mass implies a black hole superradiance window at \(M_{\rm BH} \simeq \hbar c / (G m_a) \approx 9.36 \times 10^{24}~{\rm kg}\).

\section{Quark sector (\(T^2/\mathbb{Z}_3\), \(A_4\))}
\subsection{Matrix structure and results}
\begin{equation}
M_u = \begin{pmatrix}
a_u Y_1^{(2)} & b_u Y_3^{(4)} & c_u Y_2^{(6)} \\
a_u Y_3^{(2)} & b_u Y_2^{(4)} & c_u Y_1^{(6)} \\
a_u Y_2^{(2)} & b_u Y_1^{(4)} & c_u Y_3^{(6)}
\end{pmatrix},
\end{equation}
\begin{equation}
M_d = \begin{pmatrix}
a_1 Y_1^{(6)} & b Y_3^{(2)} & c Y_2^{(2)} \\
a_1 Y_3^{(6)} & b Y_2^{(2)} & c Y_1^{(2)} \\
a_2 Y_2^{(2)} & b Y_1^{(2)} & c Y_3^{(4)}
\end{pmatrix}.
\end{equation}
Independent optimization yields \(\chi^2 < 10^{-10}\) (numerical precision limit) at the moduli given in the abstract~\cite{PDG2024}. The overall scale is unphysical; we fix \(a_u = 1\).

\section{Lepton sector (\(T^2/\mathbb{Z}_4\), \(S_4'\))}
\label{sec:lepton}
The neutrino mass matrix (Majorana) is built from the contractions of the \(\mathbf{2}\) and \(\mathbf{3}'\) representations of \(S_4'\):
\begin{equation}
M_\nu = c_1 M_{\mathbf{2}}(\tau_l) + c_2 M_{\mathbf{3}'}(\tau_l) + c_3 \mathbb{I}.
\end{equation}

\subsection{Results}
We fit the model to the three PMNS angles and the two mass-squared differences (NuFIT 5.2, normal ordering; see \url{http://www.nu-fit.org})~\cite{NuFIT2020}. The global optimization yields \(\chi^2 \approx 0.23\) at:
\begin{equation}
\tau_l = 1.5398 + 1.5511\,i, \quad c_1 = 6.36, \quad c_2 = 17.43, \quad c_3 = -17.17.
\end{equation}
The model predicts:
\begin{align}
\theta_{12} &= 33.50^\circ, \quad \theta_{13} = 8.53^\circ, \quad \theta_{23} = 49.00^\circ, \\
\Delta m^2_{21} &= 7.45 \times 10^{-5}~\text{eV}^2, \quad \Delta m^2_{31} = 2.517 \times 10^{-3}~\text{eV}^2.
\end{align}

\paragraph{Type-I See-Saw and neutrino masses.}
We fix the overall neutrino mass scale via the Type-I See-Saw mechanism: \(\Lambda_\nu = y_\nu^2 v^2 / M_6 \approx 1.78\) meV. To extract the physical masses, we perform a Takagi diagonalization of the complex symmetric Majorana matrix, \(M_\nu = U^* \mathrm{diag}(m_i) U^\dagger\), which is equivalent to diagonalizing \(M_\nu M_\nu^\dagger\). We find the dimensionless eigenvalues \(w_0 \approx 13.12\), \(w_1 \approx 36.43\), and \(w_2 \approx 800.82\). The physical masses are given by \(m_i = \sqrt{w_i} \Lambda_\nu\), which yields:
\begin{align}
m_1 &\approx 6.46~\text{meV}, \\
m_2 &\approx 10.77~\text{meV}, \\
m_3 &\approx 50.48~\text{meV}.
\end{align}
This corresponds to \(\Delta m^2_{21} \approx 7.45 \times 10^{-5}~\text{eV}^2\) and \(\Delta m^2_{31} \approx 2.517 \times 10^{-3}~\text{eV}^2\), in excellent agreement with NuFIT 5.2. The sum of neutrino masses is predicted to be \(\sum m_i \approx 0.068\) eV, which is well within the current Planck bound (\(\sum m_i < 0.12\) eV)~\cite{Planck2018} and will be testable by upcoming cosmological surveys.

\paragraph{Prediction for CP violation.}
The model predicts a non-zero Jarlskog invariant \(J_{cp} \approx 0.0156\). Using the standard PDG parameterization, we extract \(\sin\delta_{CP} \approx 0.47 \implies \delta_{CP} \approx 28^\circ\) or \(152^\circ\). This provides a concrete test for DUNE and Hyper-Kamiokande~\cite{DUNE2024}.

\paragraph{Neutrinoless double beta decay.}
Using the Takagi phases and the predicted neutrino masses, we calculate the effective Majorana mass:
\begin{equation}
m_{\beta\beta} = \left| \sum_i U_{ei}^2 m_i \right| \approx 8.7~\text{meV}.
\end{equation}
This value is well below the current KamLAND-Zen limit~\cite{KamLANDZen2023} and lies precisely within the target sensitivity of next-generation experiments such as LEGEND-1000 and nEXO~\cite{LEGEND2024}. This provides a highly testable prediction of the model.

\section{Baryon asymmetry via leptogenesis}
\label{sec:leptogenesis}
The model contains heavy right-handed neutrinos with mass \(M_R \sim M_6 \approx 3 \times 10^{16}\) GeV (Section~\ref{sec:lepton}), which can generate the observed baryon asymmetry through the leptogenesis mechanism~\cite{Fukugita1986}. 

Using the standard Davidson-Ibarra formula for the CP asymmetry in heavy neutrino decays~\cite{DavidsonIbarra2002}, we obtain a formal upper bound:
\begin{equation}
\varepsilon_{CP} \simeq \frac{3}{16\pi} \frac{M_1 m_3}{v^2} \sin\delta_{CP} \approx 0.69.
\end{equation}
We emphasize that this value is a formal upper bound, close to the perturbativity limit; a full calculation including flavour effects and the exact heavy neutrino mass spectrum will reduce it. The washout parameter is \(K \simeq 0.83\), giving an efficiency \(\kappa \simeq 1\). The resulting baryon asymmetry is:
\begin{equation}
\eta_B \approx \frac{28}{79} \frac{\varepsilon_{CP} \kappa}{g_*} \approx 2.3 \times 10^{-3}.
\end{equation}
This is about \(3.7 \times 10^6\) times larger than the observed value \(\eta_B^{\rm obs} \approx 6.1 \times 10^{-10}\). 

This apparent overproduction is naturally resolved by \emph{entropy dilution} from the late decay of moduli, which is a standard feature of KKLT-like stabilization~\cite{KKLT2003} (Section~\ref{sec:flux_stabilization}). We treat the dilution factor \(\Delta\) as a free parameter, fitted to match the observed value $\eta_B^{\rm obs}$. This yields \(\Delta \approx 3.7 \times 10^6\), corresponding to a moduli decay temperature \(T_{\rm decay} \sim 10^9\) GeV, which can be probed by future gravitational wave observatories. We stress that \(\eta_B\) is \emph{not} an independent prediction: it is obtained only after fitting \(\Delta\) to the observed value.

\section{Full balance}
\label{sec:balance}
\begin{table}[H]
\centering
\caption{Free parameters. The model has 19 effective parameters (17 fitted + 2 UV inputs) fitting 16 observables.}
\begin{tabular}{llc}
\toprule
Category & Parameters & Count \\
\midrule
Quark moduli & \(\tau_u, \tau_d\) & 4 \\
Quark coefficients & \(a_u, b_u, c_u, a_1, a_2, b, c\) & 7 \\
\quad fixed by calibration & \(a_u \equiv 1\) & \(-1\) \\
Lepton modulus & \(\tau_l\) & 2 \\
Lepton coefficients & \(c_1, c_2, c_3\) & 3 \\
ALP misalignment angle & \(\theta_i\) & 1 \\
Entropy dilution factor & \(\Delta\) & 1 \\
See-Saw Yukawa coupling & \(y_\nu\) & 1 \\
See-Saw geometric factor & \(\mathcal{I}\) & 1 \\
\midrule
\textbf{Total free parameters} & (17 fitted + 2 UV inputs) & \textbf{19} \\
\bottomrule
\end{tabular}
\end{table}

\paragraph{Information content.}
The model has 19 effective parameters. Of these, 17 are fitted to the 16 observables (9 quarks + 5 leptons + \(\Omega_a h^2\) + \(\eta_B\)): the 15 quark/lepton parameters, \(\theta_i\) for \(\Omega_a h^2\), and \(\Delta\) for \(\eta_B\). The remaining 2 parameters (\(y_\nu\) and \(\mathcal{I}\)) are UV inputs counted conservatively as free. The model makes 3 additional independent predictions that were not used in the fit:
(i) the Dirac CP phase \(\delta_{CP} \approx 28^\circ\) (testable at DUNE/Hyper-K),
(ii) the sum of neutrino masses \(\sum m_i \approx 0.068\) eV (testable at Euclid/DESI),
(iii) the effective Majorana mass \(m_{\beta\beta} \approx 8.7\) meV (testable at LEGEND-1000/nEXO).
Counting these predictions as falsifiable constraints yields:
\begin{equation}
N_{\rm dof}^{\rm fit} = 16 - 19 = -3, \qquad
N_{\rm dof}^{\rm total} = (16+3) - 19 = 0.
\end{equation}
The model is therefore \emph{exactly determined} when all predictions are counted, not overdetermined. A future UV derivation of \(y_\nu\) and \(\mathcal{I}\) would reduce \(N_{\rm eff}\) to 17 and yield \(N_{\rm dof}^{\rm total}=+2\).

\section{Geometric interpretation of Yukawa coefficients}
The hierarchical pattern of quark coefficients is interpreted via localization at the fixed points of \(T^2/\mathbb{Z}_3\). The Yukawa couplings are \(c_i \sim A_i \, e^{- \alpha \, d^2_{jk}}\), where \(d_{12}^2 = (1 + 2\mathrm{Re}\tau + |\tau|^2)/(9\,\mathrm{Im}\tau)\) and \(d_{13}^2 = 4\,d_{12}^2\). This explains the hierarchy with \(\alpha \sim 10\).

\section{Flux stabilization of the moduli}
\label{sec:flux_stabilization}
Following the GVW framework~\cite{GVW2000}, \(W(\tau) = a_0 + a_1 \tau + a_2 \tau^2\). Moduli stabilization requires \(D_\tau W = 0\). We find that the values \(\tau_u \approx -0.158 + 2.385i\) and \(\tau_d \approx 0.369 + 1.541i\) lie near the supersymmetric points of the fluxes \((-17, 0, 1)\) and \((-7, -1, 1)\), yielding small residuals \(|D_\tau W| \approx 0.16\) and \(0.14\). Exact stabilization requires non-perturbative terms \(A e^{-a \tau}\) (KKLT), which will shift the minimum to zero without significantly altering the moduli values.

\section{Open problems and Roadmap}
\begin{enumerate}[leftmargin=1.5em]
\item \textbf{External UV inputs.} The neutrino Yukawa coupling $y_\nu = 0.934$ and the geometric factor $\mathcal{I} = 0.3653$ are counted as free parameters in the balance of Section~10, giving $N_{\rm eff}=19$ and $N_{\rm dof}^{\rm total}=0$. A full UV completion deriving these values would reduce $N_{\rm eff}$ to 17 and restore $N_{\rm dof}^{\rm total}=+2$. We leave this derivation for future work.
\item \textbf{LOO instability.} The Leave-One-Out test shows extreme instability (e.g., mean \(|V_{ub}|\) deviation of \(+16728\%\)), typical for models with \(N_{\rm dof} < 0\). This is a known limitation.
\item \textbf{Status of \(\alpha_G\) and \(\varepsilon\).} Both remain postulates.
\item \textbf{Flux stabilization.} The current residuals \(|D_\tau W| \approx 0.16\) are near-SUSY points. A full KKLT analysis with Hessian is left for future work.
\item \textbf{Geometric origin of the lepton coefficients.} The coefficients \(c_1, c_2, c_3\) are treated as phenomenological parameters. They likely arise from overlap integrals of zero-mode wavefunctions on \(T^2/\mathbb{Z}_4\) with magnetic fluxes. Deriving them explicitly is a promising direction for future work.
\item \textbf{Detailed Boltzmann equations.} The leptogenesis estimate presented here is a leading-order calculation. A full numerical solution of the Boltzmann equations is left for future work.
\end{enumerate}

\section{Conclusions}
We have constructed a mixed-orbifold model \(T^2/\mathbb{Z}_3 \times T^2/\mathbb{Z}_4\).
\begin{itemize}
\item Quark sector: \(\chi^2 < 10^{-10}\) (numerical precision limit).
\item Lepton sector: using \(S_4'\) modular forms, \(\chi^2 \approx 0.23\), perfectly reproducing PMNS angles and mass differences.
\item Neutrino masses are fixed by See-Saw, predicting \(\sum m_i \approx 0.068\) eV.
\item CP violation: \(\delta_{CP} \approx 28^\circ\) (or \(152^\circ\)), testable at DUNE/Hyper-K.
\item ALP dark matter: \(m_a = 28.4~\mu\)eV, testable at ADMX/HAYSTAC.
\item Neutrinoless double beta decay: \(m_{\beta\beta} \approx 8.7\) meV, testable at LEGEND-1000/nEXO.
\item Baryon asymmetry: explained via leptogenesis, requiring an entropy dilution factor of \(\sim 3.7 \times 10^6\) from late moduli decays.
\end{itemize}
The model is phenomenological but highly successful, providing multiple concrete predictions for dark matter, neutrino masses, and CP violation.

% ============================================================================
% Order per NPB/passport: Conclusions → Data Availability → Declaration
% → Appendix → Bibliography. (Confirmed in v15.11.)
% ============================================================================

\section*{Data Availability}
The datasets and the Python verification script (v15.11) are available in the Zenodo repository at \url{https://doi.org/10.5281/zenodo.23582419} and in the GitHub repository at \url{https://github.com/kolbasa1972/MEML-75-Mixed-Geometry}.

\section*{Declaration of Generative AI and AI-assisted technologies in the manuscript preparation process}
During the preparation of this work, the author used AI-assisted tools to verify the Python code, refine the theoretical text, and format the LaTeX document. The author reviewed and edited the content as needed and takes full responsibility for the content of the published article.

\appendix
\newpage

\section{\(A_4\) modular forms: full q-expansions}
Definition: \(q = e^{2\pi i\tau}\).
\begin{align}
Y_1^{(2)} &= 1 + \sum_{n=1}^{17} C_1^{(n)} q^n, \\
Y_2^{(2)} &= -6\, q^{1/3}\left(1 + \sum_{n=1}^{11} C_2^{(n)} q^n\right), \\
Y_3^{(2)} &= -18\, q^{2/3}\left(1 + \sum_{n=1}^{10} C_3^{(n)} q^n\right),
\end{align}
\(C_1 = [12, 36, 12, 84, 72, 36, 96, 180, 12, 216, 144, 36, 156, 96, 84, 240, 36]\);
\(C_2 = [1, 7, 8, 18, 14, 30, 20, 42, 30, 56, 42]\);
\(C_3 = [1, 2, 5, 6, 12, 8, 14, 10, 22, 16]\).

\textbf{Convolutions:}
\begin{align}
Y_1^{(4)} &= (Y_1^{(2)})^2 - Y_2^{(2)} Y_3^{(2)}, \\
Y_2^{(4)} &= (Y_3^{(2)})^2 - Y_1^{(2)} Y_2^{(2)}, \\
Y_3^{(4)} &= (Y_2^{(2)})^2 - Y_1^{(2)} Y_3^{(2)},
\end{align}
\begin{align}
Y_1^{(6)} &= \tfrac13\left[2 Y_1^{(2)} Y_1^{(4)} - Y_2^{(2)} Y_3^{(4)} - Y_3^{(2)} Y_2^{(4)}\right], \\
Y_2^{(6)} &= \tfrac13\left[2 Y_3^{(2)} Y_3^{(4)} - Y_1^{(2)} Y_2^{(4)} - Y_2^{(2)} Y_1^{(4)}\right], \\
Y_3^{(6)} &= \tfrac13\left[2 Y_2^{(2)} Y_2^{(4)} - Y_3^{(2)} Y_1^{(4)} - Y_1^{(2)} Y_3^{(4)}\right].
\end{align}

\section{Numerical fit procedure}
The global fit was performed using \texttt{differential\_evolution} (SciPy):
4 seeds (42, 43, 44, 45); maxiter$=800$; popsize$=25$; tol$=10^{-8}$.

\emph{Quark sector.} Bounds:
\(\mathrm{Re}\,\tau_{u,d} \in [-0.5, 0.5]\),
\(\mathrm{Im}\,\tau_u \in [0.87, 2.5]\),
\(\mathrm{Im}\,\tau_d \in [0.87, 1.60]\),
\(c_i \in [10^{-8}, 10^2]\).

\section{Reproduction of \(\tau\) and \(\chi^2\)}
\begin{table}[H]
\centering
\caption{Independent DE runs (quark sector).}
\begin{tabular}{lcccc}
\toprule
Source & \(\mathrm{Re}\,\tau_u\) & \(\mathrm{Im}\,\tau_u\) & \(\mathrm{Re}\,\tau_d\) & \(\mathrm{Im}\,\tau_d\) \\
\midrule
Run 1 & \(-0.15779\) & \(2.38518\) & \(+0.36939\) & \(1.54105\) \\
Run 2 & \(-0.15779\) & \(2.38518\) & \(+0.36939\) & \(1.54105\) \\
\bottomrule
\end{tabular}
\end{table}
The fitted moduli \(\mathrm{Im}\,\tau\) are reproducible across seeds to \(10^{-3}\). The quantity \(\chi^2\) at the minimum is a machine zero; we quote it as \(\chi^2 < 10^{-10}\).

\section{Proof of coefficient degeneracy}
At fixed \(\tau\), multiply all \(c_i\) by \(k\):
\begin{table}[H]
\centering
\caption{Invariance of \(\chi^2\) under \(c_i \to k c_i\).}
\begin{tabular}{cc}
\toprule
\(k\) & \(\chi^2\) \\
\midrule
\(0.5\)  & \(<10^{-15}\) \\
\(1.0\)  & \(<10^{-15}\) \\
\(1.87\) & \(<10^{-15}\) \\
\(10.0\) & \(<10^{-15}\) \\
\bottomrule
\end{tabular}
\end{table}
The residual values differ only at the level of numerical noise; the invariance is exact up to machine precision.

\section{Flatness of \(\mathrm{Re}\,\tau\)}
\paragraph{Procedure.}
The flatness scan below was obtained by refitting the remaining parameters at each value of \(\mathrm{Re}\,\tau\). The corresponding refit cycle is computationally expensive and is not part of the fast verification script.

Scan in \(\mathrm{Re}\,\tau_u\): width of \(\chi^2<10\): \([-0.40, +0.00]\) (\(0.40\)).
Scan in \(\mathrm{Re}\,\tau_d\): width of \(\chi^2<10\): \([+0.30, +0.50]\) (\(0.20\)).

\begin{thebibliography}{99}
\bibitem{Planck2018} Planck Collaboration, \emph{Planck 2018 results. VI. Cosmological parameters}, Astron. Astrophys. \textbf{641}, A6 (2020), arXiv:1807.06209, DOI:~\href{https://doi.org/10.1051/0004-6361/201833910}{10.1051/0004-6361/201833910}.
\bibitem{PDG2024} R.~L. Workman et al. (Particle Data Group), \emph{Review of Particle Physics}, Prog. Theor. Exp. Phys. \textbf{2022}, 083C01 (2022), DOI:~\href{https://doi.org/10.1093/ptep/ptac097}{10.1093/ptep/ptac097}.
\bibitem{ADMX2018} ADMX Collaboration, Phys. Rev. Lett. \textbf{120}, 151301 (2018), arXiv:1804.05750.
\bibitem{HAYSTAC2021} HAYSTAC Collaboration, \emph{A quantum-enhanced search for dark matter axions}, Nature \textbf{590}, 238 (2021), arXiv:2008.01853.
\bibitem{Feruglio2017} F.~Feruglio, \emph{Are neutrino masses modular forms?}, Eur. Phys. J. C \textbf{77}, 713 (2017), arXiv:1706.08749.
\bibitem{Kobayashi2019} T.~Kobayashi, Y.~Shimizu, K.~Takagi, M.~Tanimoto, T.~H. Tatsuishi, \emph{Modular \(S_3\)-invariant flavor model in SU(5) grand unified theory}, PTEP \textbf{2020}, 053B05 (2020), arXiv:1906.10341.
\bibitem{Novichkov2019} P.~P. Novichkov, J.~T. Penedo, S.~T. Petcov, A.~V. Titov, JHEP \textbf{04} (2019) 005, arXiv:1811.04933.
\bibitem{Ding2021} G.~J. Ding, S.~F. King, C.~Y. Yao, \emph{Modular \(S_4 \times SU(5)\) GUT}, Phys. Rev. D \textbf{104}, 055034 (2021), arXiv:2103.16311.
\bibitem{Nilles2020} H.~P. Nilles, S.~Ramos-Sánchez, P.~K.~S. Vaudrevange, \emph{Lessons from eclectic flavor symmetries}, Nucl. Phys. B \textbf{957} (2020) 115098, arXiv:2004.05200.
\bibitem{Baur2024} A.~Baur, H.~P. Nilles, S.~Ramos-Sánchez, A.~Trautner, P.~K.~S. Vaudrevange, \emph{The eclectic flavor symmetries of \(\mathbb{T}^2/\mathbb{Z}_K\) orbifolds}, JHEP \textbf{09} (2024) 159, arXiv:2405.20378.
\bibitem{Fukugita1986} M.~Fukugita, T.~Yanagida, \emph{Baryogenesis without Grand Unification}, Phys. Lett. B \textbf{174} (1986) 45, DOI:~\href{https://doi.org/10.1016/0370-2693(86)91126-3}{10.1016/0370-2693(86)91126-3}.
\bibitem{NuFIT2020} I.~Esteban, M.~C. Gonzalez-Garcia, M.~Maltoni, T.~Schwetz, \emph{The fate of hints: updated global analysis of three-flavor neutrino oscillations}, JHEP \textbf{09} (2020) 178, arXiv:2007.14792 (NuFIT 5.2, 2022).
\bibitem{KamLANDZen2023} KamLAND-Zen Collaboration, Phys. Rev. Lett. \textbf{130}, 051801 (2023), arXiv:2203.02139.
\bibitem{DUNE2024} DUNE Collaboration, \emph{DUNE Far Detector Technical Design Report, Volume II: Physics Program}, arXiv:2002.03005.
\bibitem{LEGEND2024} LEGEND Collaboration, \emph{LEGEND-1000 Conceptual Design Report}, arXiv:2107.11462.
\bibitem{DavidsonIbarra2002} S.~Davidson, A.~Ibarra, \emph{A lower bound on the right-handed neutrino mass from leptogenesis}, Phys. Lett. B \textbf{535} (2002) 25, arXiv:hep-ph/0202239.
\bibitem{KKLT2003} S.~Kachru, R.~Kallosh, A.~Linde, S.~P. Trivedi, \emph{de Sitter vacua in string theory}, Phys. Rev. D \textbf{68}, 046005 (2003), arXiv:hep-th/0301240.
\bibitem{GVW2000} S.~Gukov, C.~Vafa, E.~Witten, \emph{CFT's from Calabi-Yau four-folds}, Nucl. Phys. B \textbf{584} (2000) 69, arXiv:hep-th/9906070.
\end{thebibliography}

\end{document}
